\documentclass[12pt]{report}

\usepackage[margin=1in]{geometry}
\usepackage{setspace}
\usepackage{enumitem}
\usepackage{titlesec}
\usepackage{xcolor}
\usepackage{pdflscape}
\usepackage{graphicx}

\definecolor{background_color}{RGB}{27,25,25}
\pagecolor{background_color}
\color{white}
\titleformat{\section}{\large\bfseries}{\thesection}{1em}{}
\setstretch{1.2}

\renewcommand{\thesection}{HW \#\arabic{section}}   % sections: HW #1, HW #2, ...
\renewcommand{\thesubsection}{\arabic{section}.\arabic{subsection}}   % subsections: 1.1, 1.2, ...
\titleformat{\section}
  {\normalfont\Large\bfseries\clearpage}  % \textless-- page break added here
  {HW \#\arabic{section}}
  {1em}
  {}
  
\title{UMBC Parking Management System\\Database Design Proposal}
\author{Mateo Jacome, Jason Chen}
\date{\today}

\begin{document}

\maketitle

\section{Executive Summary}

\subsection{Problem Statement}

Parking at a large commuter campus such as UMBC is inefficient, stressful, and difficult to enforce. Students frequently waste valuable time searching for open spaces. Visitors struggle to locate valid parking without clear guidance. Faculty and staff require access to restricted or premium areas. Meanwhile, enforcement is largely reactive and labor-intensive, relying on officers instead of automated validation mechanisms.

The core problem is the absence of a centralized, real-time database system that integrates parking permits, lot and spot management, vehicle data, reservations, sensor occupancy detection, and automated enforcement into a unified platform. Without such integration, parking operations lack efficiency, transparency, and scalability.

The proposed UMBC Parking Management System addresses this issue by implementing a PostgreSQL-backed database that manages parking as a real-time, data-driven ecosystem.

\subsection{Parking Users}

The system supports multiple categories of users:

\begin{itemize}
    \item \textbf{Commuter Students} -- Purchase student permits and park in designated student lots.
    \item \textbf{Faculty and Staff} -- Hold faculty/staff permits and access restricted premium zones.
    \item \textbf{Visitors and Guest Lecturers} -- Reserve temporary parking spaces online and pay per visit.
\end{itemize}

\subsection{Scope of Implementation}

This project will implement a PostgreSQL database system that includes:

\begin{enumerate}
    \item \textbf{User and Vehicle Management}
    \begin{itemize}
        \item Store user profiles (students, faculty, visitors).
        \item Associate multiple vehicles with individual users.
        \item Store license plate information for validation.
    \end{itemize}

    \item \textbf{Permit Management}
    \begin{itemize}
        \item Define permit types (student, faculty, visitor).
        \item Assign permits to users.
        \item Enforce expiration dates.
        \item Restrict lot and zone access by permit type.
    \end{itemize}

    \item \textbf{Lot and Spot Management}
    \begin{itemize}
        \item Define parking lots, zones, rows, and individual spots.
        \item Track real-time occupancy per spot.
        \item Associate spots with allowed permit types.
    \end{itemize}

    \item \textbf{Reservation System}
    \begin{itemize}
        \item Allow visitors to reserve specific spots.
        \item Prevent double-booking through constraints and transactions.
        \item Lock spots upon confirmed payment.
    \end{itemize}

    \item \textbf{Sensor Integration}
    \begin{itemize}
        \item Update occupancy when a ground sensor detects a vehicle.
        \item Validate vehicle authorization against zone restrictions.
    \end{itemize}

    \item \textbf{Automated Ticketing}
    \begin{itemize}
        \item Generate violations for unauthorized parking.
        \item Store fine amounts and violation details.
        \item Record notification events.
    \end{itemize}

    \item \textbf{Payments}
    \begin{itemize}
        \item Record ticket payments.
        \item Track outstanding balances.
    \end{itemize}

    \item \textbf{Reporting and Analytics}
    \begin{itemize}
        \item Lot utilization rates.
        \item Revenue reports.
        \item Violation frequency reports.
        \item Active and expired permit tracking.
        \item Peak usage analysis.
    \end{itemize}

    \item \textbf{Concurrency Control}
    \begin{itemize}
        \item Prevent race conditions during reservation.
    \end{itemize}
\end{enumerate}

This scope focuses specifically on the database design, implementation, constraints, indexing, realistic operations, and concurrency control. It does not include a full mobile or web application interface.

\subsection{Concrete Use Cases}

The database must support the following operations:

\begin{itemize}
    \item Register a new user (student, faculty, or visitor).
    \item Add one or more vehicles to a user account.
    \item Purchase and assign a parking permit to a user.
    \item Define parking lots and individual parking spots.
    \item Associate permit types with allowed zones.
    \item Query real-time available parking spots in a specific lot.
    \item Create a visitor reservation for a specific date and time.
    \item Prevent double-booking of a parking spot.
    \item Update occupancy when a sensor detects a vehicle.
    \item Validate a vehicle's permit against the spot's allowed zone.
    \item Automatically generate a ticket for unauthorized parking.
    \item Record ticket payments and mark violations as resolved.
    \item Generate reports on lot occupancy over time.
    \item Generate revenue reports from permits, reservations, and fines.
    \item Detect expired permits and flag vehicles as invalid.
    \item Handle simultaneous reservation attempts without data corruption.
\end{itemize}

\section{Requirements \& Conceptual Design}
This section will more rigourously define the requirements, constraints, roles, as well as include an ER diagram.

\subsection{ER Diagram}
    To see full image click on \verb|er_diagram.png|
    \begin{center}
    \includegraphics[width=\textwidth, keepaspectratio]{er_diagram.png}
    \end{center}

\subsection{Requirements and Constraints}

The following business rules define constraints for the UMBC Parking Management System. Each rule is labeled with how it is enforced: (A) through database table constraints such as PRIMARY KEY, FOREIGN KEY, UNIQUE, or CHECK, or (B) through triggers, functions, or application logic.

\begin{enumerate}[label=\textbf{Rule \arabic*:}]

\item Each user may register multiple vehicles, but each license plate must be unique across the system.

Enforcement: \textbf{(A)} A UNIQUE constraint is applied to the car\_info.plate column.

\item Each parking spot within a lot must be uniquely identified by its lot and space number.

Enforcement: \textbf{(A)} A composite UNIQUE constraint is applied to (lot\_id, space) in the parking\_spot table.

\item A reservation must have a start time that occurs before its end time.

Enforcement: \textbf{(A)} A CHECK constraint ensures start\_time \textless end\_time.

\item A permit must expire after it is issued.

Enforcement: \textbf{(A)} A CHECK constraint ensures expiration\_date \textgreater issued\_date.

\item Monetary values for fines and payments must be positive.

Enforcement: \textbf{(A)} CHECK constraints ensure values such as fine\_amount\_cents \textgreater 0 and amount\_cents \textgreater 0.

\item A parking spot cannot be reserved by two users for overlapping time intervals.

Enforcement: \textbf{(B)} Implemented using a PostgreSQL trigger that looks at the time range derived from start\_time and end\_time.

\item Only vehicles with a valid and unexpired permit may park in permit-restricted lots.

Enforcement: \textbf{(B)} Implemented using a trigger or application-level validation that checks the vehicle's permit type and expiration date.

\item Visitor reservations must be confirmed before the space is considered locked.

Enforcement: \textbf{(B)} Application logic ensures reservations transition from PENDING to CONFIRMED only after successful payment.

\item Sensor events must record an arrival or departure event with a timestamp.

Enforcement: \textbf{(A)} Implemented using the sensor\_event\_type enum and a NOT NULL timestamp constraint.

\item If a vehicle parks without authorization (no valid permit or reservation), the system must generate a parking violation ticket.

Enforcement: \textbf{(B)} Implemented using a trigger that listens for arrival events and inserts a record into the tickets table when authorization fails.

\end{enumerate}

\subsection{Concurrency Scenario}

A concurrency scenario occurs when two users attempt to reserve the same parking spot at nearly the same time. For example, User A and User B both attempt to reserve parking spot S in Lot 3 from 9:00 AM to 12:00 PM.

Without proper concurrency control, both reservations could be inserted into the database, creating a double booking. The system must guarantee that only one reservation succeeds.

To enforce this requirement, the database uses a PostgreSQL trigger that runs BEFORE insertion into the reservations table, that prevents overlapping time ranges for the same parking spot. When the first reservation transaction commits successfully, any subsequent reservation that overlaps with that time interval will automatically fail.

This approach ensures correctness even when multiple transactions occur simultaneously.

\subsection{Rationale}

The design separates long-term policy information from transactional records in order to improve flexibility and maintainability. Permit classifications (such as student, faculty, or visitor) represent policy rules that determine which parking areas are accessible to different groups. By separating permit types from individual permit records, the system allows parking policies to evolve independently from historical permit data.

Reservations are modeled using explicit start and end timestamps rather than a single date field so that the system can support flexible time windows. This design enables features such as hourly reservations, overlapping time checks, and historical analysis of parking usage patterns.

Sensor events are stored as individual event records instead of directly updating a parking spot's status. This event-driven model preserves a complete history of arrivals and departures, which is valuable for analytics, auditing, and debugging sensor behavior. A trigger or background process can interpret these events to update occupancy status, validate permits, or generate tickets when violations occur.

Finally, the schema models relationships such as vehicles belonging to users and reservations referencing parking spots using foreign keys. These relationships ensure referential integrity and make it possible to query the system efficiently for operations such as finding available parking, verifying permits, or generating enforcement reports.



\end{document}